\documentclass[11pt,a4paper]{article}

%======================================================================
% PACKAGES
%======================================================================
\usepackage[utf8]{inputenc}
\usepackage[T1]{fontenc}
\usepackage[margin=2.3cm]{geometry}
\usepackage{amsmath,amssymb}
\usepackage{listings}
\usepackage{xcolor}
\usepackage{hyperref}
\usepackage{enumitem}
\usepackage{fancyhdr}
\usepackage{titlesec}
\usepackage{tcolorbox}
\usepackage{graphicx}
\usepackage{array}
\usepackage{booktabs}

%======================================================================
% STYLE
%======================================================================
\definecolor{codebg}{rgb}{0.96,0.96,0.92}
\definecolor{codekw}{rgb}{0.10,0.10,0.65}
\definecolor{codestr}{rgb}{0.55,0.10,0.10}
\definecolor{codecom}{rgb}{0.20,0.45,0.20}
\definecolor{boxblue}{rgb}{0.85,0.92,1.0}
\definecolor{boxyellow}{rgb}{1.0,0.97,0.80}
\definecolor{boxred}{rgb}{1.0,0.88,0.88}

\lstdefinestyle{C}{
    language=C,
    backgroundcolor=\color{codebg},
    basicstyle=\ttfamily\footnotesize,
    keywordstyle=\color{codekw}\bfseries,
    stringstyle=\color{codestr},
    commentstyle=\color{codecom}\itshape,
    breaklines=true,
    showstringspaces=false,
    frame=single,
    rulecolor=\color{black!30},
    numbers=left,
    numberstyle=\tiny\color{gray},
    captionpos=b,
    tabsize=4
}
\lstdefinestyle{shell}{
    backgroundcolor=\color{codebg},
    basicstyle=\ttfamily\footnotesize,
    breaklines=true,
    frame=single,
    rulecolor=\color{black!30}
}

\newtcolorbox{keyidea}{colback=boxblue, colframe=blue!50!black, title=Key Idea}
\newtcolorbox{warning}{colback=boxred, colframe=red!50!black, title=Watch Out}
\newtcolorbox{note}{colback=boxyellow, colframe=orange!60!black, title=Note}

\titleformat{\section}{\Large\bfseries\color{blue!40!black}}{\thesection}{1em}{}
\titleformat{\subsection}{\large\bfseries\color{blue!30!black}}{\thesubsection}{1em}{}

\setlength{\headheight}{14pt}
\pagestyle{fancy}
\fancyhf{}
\fancyhead[L]{Firmware Design --- In-Depth Notes}
\fancyhead[R]{\thepage}

\hypersetup{
    colorlinks=true,
    linkcolor=blue!50!black,
    urlcolor=blue!60!black,
}

\setlength{\parskip}{6pt}
\setlength{\parindent}{0pt}

%======================================================================
\title{\textbf{Firmware Design --- The Big "What the Hell is Going On" Guide}\\
       \large An in-depth walkthrough of the course syllabus}
\author{Course companion notes}
\date{May 2026}

\begin{document}
\maketitle
\tableofcontents
\newpage

%======================================================================
\section{Classic vs.\ Embedded Systems}
%======================================================================

A ``classic'' data processing system is what you are reading these notes on:
a PC or a server. An \textbf{embedded system} is a computer that you usually
don't \emph{see} as a computer: the thing inside your washing machine, car ECU,
drone flight controller, smart thermostat, or pacemaker. Both run code, but
they live in very different worlds.

\subsection{Hardware differences}

\begin{tabular}{p{0.27\linewidth} p{0.34\linewidth} p{0.34\linewidth}}
\toprule
\textbf{Aspect} & \textbf{Classic (PC/Server)} & \textbf{Embedded (MCU)} \\
\midrule
CPU & GHz, multi-core, out-of-order, deep pipeline, branch predictor, big caches & MHz to few hundred MHz, often single-core (Cortex-M, AVR), shallow pipeline, little or no cache \\
Memory & GB of DRAM, virtual memory via MMU & KB to a few MB of \emph{on-chip} SRAM; no MMU (usually an MPU at most) \\
Storage & TB of SSD/HDD with a file system & KB--MB of internal Flash; sometimes external SPI Flash; no FS unless you bolt one on \\
Peripherals & PCIe, USB, NIC; abstracted by OS drivers & GPIO, UART, SPI, I\textsuperscript{2}C, ADC, DAC, timers, CAN; you talk to them register-by-register \\
Power & 10s--100s of W, fans, heatsinks & mA or µA; often battery powered, sometimes energy harvesting \\
Cost & 100s--1000s of \$, depreciates over years & cents to a few \$, sold by the million \\
\bottomrule
\end{tabular}

\subsection{Software differences}

\begin{itemize}[leftmargin=*]
    \item \textbf{Operating system.} A PC has Linux/Windows/macOS with a kernel,
    drivers, multitasking, virtual memory, security boundaries, paging, etc.
    An MCU often runs \emph{bare-metal} (your code \emph{is} everything), or at
    most a small RTOS (FreeRTOS, Zephyr, ThreadX) that gives you tasks,
    queues and mutexes but no MMU and no protected processes.
    \item \textbf{Real-time constraints.} A PC's job is throughput
    (``finish quickly on average''). An embedded system frequently has
    \emph{deadlines} (``never miss a sample window or the motor melts'').
    A 1 ms hiccup on Linux is invisible; on a brushless-motor controller it
    can be a fire.
    \item \textbf{Determinism.} You avoid features that introduce
    unpredictability: dynamic memory allocation, garbage collection,
    unbounded loops/recursion, virtual memory faults, etc.
    \item \textbf{Resource consciousness.} Every byte of Flash and every
    byte of RAM matters. You profile not just runtime but \emph{stack depth}
    and \emph{ROM usage}.
    \item \textbf{Update model.} On a PC you ``apt upgrade''. On an embedded
    device you must design a firmware update mechanism (bootloader, dual-bank
    flash, signature checking) or accept that the firmware is frozen forever.
    \item \textbf{Failure model.} On a PC, the OS catches and reports faults.
    On a bare-metal MCU, a wild pointer either silently corrupts memory or
    triggers a HardFault that, if not handled, locks the chip --- often the
    only feedback is ``the LED stopped blinking''.
\end{itemize}

\begin{keyidea}
An embedded engineer writes code that runs \emph{forever}, on hardware that
\emph{costs cents}, against \emph{real-world physics}, with \emph{no operator}
to restart the process. That changes every single engineering decision compared
to PC software.
\end{keyidea}

%======================================================================
\section{Microprocessor Interrupts Management}
%======================================================================

An \textbf{interrupt} is a hardware-triggered jump in the CPU's normal
execution flow. The CPU is doing something, an event happens (a UART byte
arrives, a timer expires, a button is pressed), and the hardware forces the
CPU to stop what it was doing, save its state, jump to a piece of code called
the \textbf{Interrupt Service Routine (ISR)}, run it, then resume the original
work as if nothing happened.

\subsection{Why interrupts at all?}

The alternative is \textbf{polling}: a loop that constantly checks each
peripheral, ``do you have data? do you have data? do you have data?''.
Polling wastes CPU, increases latency to the worst-case polling period, and
makes power management impossible (CPU can never sleep). Interrupts let the
CPU sleep or do useful work and react only when needed.

\subsection{The mechanism, step by step}

\begin{enumerate}[leftmargin=*]
    \item Peripheral asserts an IRQ line (e.g.\ ``USART2\_RX has a byte'').
    \item The Interrupt Controller (\textbf{NVIC} on ARM Cortex-M) prioritises
    pending interrupts and, if priority allows, signals the core.
    \item The core finishes the current instruction, then automatically pushes
    a context (on Cortex-M: \texttt{xPSR, PC, LR, R12, R3, R2, R1, R0}) onto
    the current stack. This is the \emph{hardware-stacked frame}.
    \item The core loads the PC from the \textbf{Interrupt Vector Table}
    entry corresponding to that IRQ number. The vector table is just an array
    of function pointers at a fixed memory address (start of Flash on Cortex-M).
    \item The ISR runs. It must clear the source of the interrupt (e.g.\
    read the data register, write the ``flag clear'' bit), otherwise the
    interrupt re-fires forever.
    \item ISR returns by executing a special return value (\texttt{EXC\_RETURN}
    on Cortex-M). Hardware unstacks the frame and resumes the interrupted code.
\end{enumerate}

\subsection{Important concepts}

\begin{description}[leftmargin=*, style=nextline]
    \item[Maskable vs.\ Non-Maskable (NMI).] Maskable IRQs can be globally
    disabled (\texttt{cpsid i} on Cortex-M, \texttt{cli} on AVR). NMI cannot
    be masked --- reserved for catastrophic things (clock failure, watchdog
    pre-warn).
    \item[Priority \& preemption.] Each interrupt has a priority. A higher-
    priority interrupt can preempt a lower-priority ISR currently running.
    Same-priority interrupts do not preempt each other; they queue.
    \item[Re-entrancy.] If you call a function from both an ISR and main
    code, it must be re-entrant: no static state, or properly protected
    with critical sections. \texttt{strtok}, \texttt{rand}, \texttt{malloc}
    are not re-entrant.
    \item[Latency.] Interrupt latency = time from event to first ISR
    instruction. Cortex-M3/M4: 12 cycles deterministic, including the
    hardware stacking. This determinism is the whole point of Cortex-M for
    real-time.
    \item[Tail-chaining \& late arrival.] Cortex-M optimisations: if a new
    IRQ is pending when an ISR finishes, the unstack/re-stack is skipped
    (tail-chaining). If a higher-priority IRQ arrives during the stacking
    of a lower one, the higher one runs first (late arrival).
\end{description}

\begin{warning}
The \#1 ISR mistake: doing too much work inside it. ISRs should be
\emph{short}. Read the byte, push it to a queue, signal a task, return.
Long ISRs increase worst-case latency for all other interrupts and make the
system look fine in testing and explode in the field.
\end{warning}

\begin{lstlisting}[style=C, caption={Typical ISR pattern (STM32 HAL)}]
volatile uint8_t rx_buf[64];
volatile uint8_t rx_head = 0;

void USART2_IRQHandler(void)
{
    if (USART2->SR & USART_SR_RXNE) {
        rx_buf[rx_head++ & 0x3F] = USART2->DR; /* reading DR clears RXNE */
    }
}
\end{lstlisting}

Notice the \texttt{volatile}: the ISR and main code share \texttt{rx\_head}
and \texttt{rx\_buf}, so the compiler must not cache these in registers
across uses.

%======================================================================
\section{Executable Code Generation --- The GCC Model}
%======================================================================

A C program does not magically become an executable. GCC orchestrates
\emph{four} distinct stages, each of which can be stopped and inspected.

\subsection{The four stages}

\begin{center}
\texttt{.c} $\xrightarrow{\text{preprocess}}$ \texttt{.i}
        $\xrightarrow{\text{compile}}$ \texttt{.s}
        $\xrightarrow{\text{assemble}}$ \texttt{.o}
        $\xrightarrow{\text{link}}$ \texttt{a.out / .elf}
\end{center}

\begin{enumerate}[leftmargin=*]
    \item \textbf{Preprocessing} (\texttt{cpp}). Expands \texttt{\#include},
    \texttt{\#define}, \texttt{\#if}, removes comments. Output is still C,
    but with all headers inlined and macros substituted. No compilation yet.
    \item \textbf{Compilation} (\texttt{cc1}). Parses the C, does optimisations,
    emits \emph{assembly} for the target architecture. This is where things
    like inlining, dead-code elimination, and register allocation happen.
    \item \textbf{Assembly} (\texttt{as}). Turns assembly mnemonics into
    machine code, producing a relocatable \textbf{object file} (\texttt{.o}).
    Object files contain code and data, plus a \emph{symbol table} and
    \emph{relocation entries} (placeholders to be patched at link time).
    \item \textbf{Linking} (\texttt{ld}). Combines all the \texttt{.o} files
    and libraries, resolves cross-references, lays out the final memory map
    according to a \emph{linker script}, and emits the executable (ELF on
    Linux / embedded, PE on Windows, Mach-O on macOS).
\end{enumerate}

\subsection{Stopping at each stage}

\begin{lstlisting}[style=shell, caption={Useful gcc invocations}]
# Preprocess only -- look at the expanded source
gcc -E hello.c -o hello.i

# Compile to assembly -- look at what gcc emitted
gcc -S hello.c -o hello.s
gcc -S -O2 hello.c -o hello_opt.s   # with optimisation

# Compile + assemble, but don't link -- get the .o
gcc -c hello.c -o hello.o

# Full pipeline
gcc hello.c -o hello

# Useful diagnostic flags
gcc -Wall -Wextra -Wpedantic -O2 -g hello.c -o hello
gcc -v hello.c          # verbose: show every sub-tool invocation
gcc -save-temps hello.c # keep .i and .s alongside the output
\end{lstlisting}

\subsection{objdump: read the binaries you produce}

\texttt{objdump} is to object files what \texttt{less} is to text files.
You should reach for it constantly when ``something is weird''.

\begin{lstlisting}[style=shell, caption={objdump cheat sheet}]
objdump -h hello.o      # list sections and their sizes
objdump -d hello.o      # disassemble code sections
objdump -D hello.o      # disassemble EVERYTHING (including data)
objdump -t hello.o      # symbol table
objdump -r hello.o      # relocation entries (what the linker still has to patch)
objdump -s hello.o      # hex+ASCII dump of all sections (good for .rodata)
objdump -x hello.o      # everything: headers, symbols, relocations
objdump -S hello        # interleave source with disasm (needs -g at compile)

# For embedded: pick the right toolchain prefix
arm-none-eabi-objdump -d firmware.elf
\end{lstlisting}

\begin{keyidea}
If you can read what \texttt{objdump -d} produces, you can debug almost any
``the C code looks right but the chip does the wrong thing'' problem.
Optimisers are aggressive --- the only ground truth is the disassembly.
\end{keyidea}

%======================================================================
\section{Cross-Compilation Environments for Embedded Systems}
%======================================================================

A native compiler runs on \emph{the same} architecture it produces code for:
gcc on x86 producing x86 binaries. A \textbf{cross-compiler} runs on one
architecture (the \emph{host}, usually x86\_64) and produces code for another
(the \emph{target}, usually ARM, RISC-V, AVR, etc.).

\subsection{Why cross-compile?}

Because the target MCU is almost certainly too small to host a C compiler:
your STM32F4 has 192 KB of RAM; GCC needs hundreds of MB. So you build on
your laptop, then transfer the resulting \texttt{.elf}/\texttt{.bin} to the
chip via a debugger probe (ST-Link, J-Link) or a bootloader (DFU, UART).

\subsection{The toolchain triplet}

A toolchain is named by a \emph{triplet} of the form
\texttt{arch-vendor-os-abi}. For embedded ARM you'll see:

\begin{itemize}[leftmargin=*]
    \item \texttt{arm-none-eabi-} --- ARM CPU, no OS (\texttt{none}),
    EABI calling convention. Used for \emph{bare-metal} and RTOS firmware on
    Cortex-M / Cortex-R. This is the ST/CubeIDE default.
    \item \texttt{arm-linux-gnueabihf-} --- ARM CPU, Linux user-space,
    glibc, hard-float ABI. Used when targeting embedded Linux boards
    (Raspberry Pi, BeagleBone).
    \item \texttt{riscv64-unknown-elf-} --- RISC-V 64, bare-metal.
\end{itemize}

Each tool is prefixed: \texttt{arm-none-eabi-gcc}, \texttt{-as}, \texttt{-ld},
\texttt{-objdump}, \texttt{-objcopy}, \texttt{-gdb}, \texttt{-size}.

\subsection{What changes vs.\ native gcc?}

\begin{itemize}[leftmargin=*]
    \item You must tell it the exact CPU and FPU:
    \texttt{-mcpu=cortex-m4 -mthumb -mfpu=fpv4-sp-d16 -mfloat-abi=hard}.
    \item There is no \texttt{libc} that talks to an OS; you use
    \textbf{newlib} or \textbf{newlib-nano}, which expects you to provide
    \emph{syscalls} (\texttt{\_write}, \texttt{\_sbrk}, \texttt{\_read}, ...).
    \item You must supply a \emph{linker script} (\texttt{.ld}) describing
    Flash and RAM regions.
    \item You must supply \emph{startup code} (\texttt{startup\_stm32xxx.s})
    that sets up the stack, zeroes \texttt{.bss}, copies \texttt{.data} from
    Flash to RAM, then calls \texttt{main()}.
    \item The output \texttt{.elf} usually has to be converted to raw binary
    or Intel HEX for the flasher:
    \texttt{arm-none-eabi-objcopy -O binary fw.elf fw.bin}.
\end{itemize}

%======================================================================
\section{Sections in Object Files and Executables}
%======================================================================

An ELF binary is divided into \textbf{sections}, each holding one ``kind'' of
content. The linker decides where each section lives in memory.

\subsection{The canonical sections}

\begin{tabular}{p{0.18\linewidth} p{0.45\linewidth} p{0.32\linewidth}}
\toprule
\textbf{Section} & \textbf{What's in it} & \textbf{Lives in (MCU)} \\
\midrule
\texttt{.text}   & Executable code (machine instructions) & Flash \\
\texttt{.rodata} & Read-only data: string literals, \texttt{const} globals, lookup tables & Flash \\
\texttt{.data}   & Initialised, writable globals/statics (e.g.\ \texttt{int x = 5;}) & RAM (loaded from Flash at boot) \\
\texttt{.bss}    & Zero-initialised globals/statics (e.g.\ \texttt{int y;}) & RAM (zeroed at boot) \\
\texttt{.stack}  & Function call frames, local variables & RAM \\
\texttt{.heap}   & Dynamically allocated blocks (\texttt{malloc}) & RAM \\
\texttt{.isr\_vector} & Interrupt vector table & Beginning of Flash \\
\bottomrule
\end{tabular}

\subsection{Why this matters in embedded}

On a PC you almost never think about sections. On an MCU you do, because:

\begin{itemize}[leftmargin=*]
    \item Flash and RAM are \emph{separate physical memories} at different
    addresses. The linker must tell each section which one it lives in.
    \item \texttt{.data} is tricky: it needs an initial value at boot,
    but it must be writable. Solution: a \emph{copy} is stored in Flash, and
    the startup code copies it to RAM before \texttt{main()}.
    \item Sections let you place specific code/data at specific addresses
    (e.g.\ DMA buffer in a particular SRAM bank --- see Section~\ref{sec:dma}).
\end{itemize}

\subsection{Custom sections via attributes}

\begin{lstlisting}[style=C, caption={Placing variables in custom sections}]
/* Put this buffer in section ".dma_buffer", defined in linker script */
__attribute__((section(".dma_buffer"))) uint8_t dma_buf[1024];

/* Put a function in RAM for speed (must be copied at boot) */
__attribute__((section(".ramfunc"))) void fast_isr(void) { ... }
\end{lstlisting}

The linker script then has matching lines like:
\begin{lstlisting}[style=shell]
.dma_buffer (NOLOAD) : { *(.dma_buffer) } > SRAM2
\end{lstlisting}

%======================================================================
\section{Memory Allocation of Static vs.\ Local Variables}
%======================================================================

This is a topic where ``intuition from PC code'' lies to you. Let's be precise.

\subsection{Static variables}

\textbf{Static} = \emph{static storage duration}. These exist for the whole
program lifetime. That includes:
\begin{itemize}[leftmargin=*]
    \item File-scope globals (\texttt{int counter;}).
    \item File-scope \texttt{static} globals (same lifetime, file-local visibility).
    \item Function-scope \texttt{static} locals (\texttt{static int n = 0;}
    inside a function --- exists across calls, only initialised once).
\end{itemize}

Where do they live?
\begin{itemize}[leftmargin=*]
    \item Initialised to nonzero $\rightarrow$ \texttt{.data} (RAM at run-time,
    initial value sitting in Flash).
    \item Initialised to zero or uninitialised $\rightarrow$ \texttt{.bss}
    (zeroed RAM, no Flash cost for initial values).
    \item \texttt{const} qualified $\rightarrow$ \texttt{.rodata} (Flash on MCU).
\end{itemize}

\subsection{Local (automatic) variables}

A function-local non-\texttt{static} variable has \emph{automatic} storage:
created on function entry, destroyed on return. Conceptually on the
\textbf{stack}. \emph{In practice, the optimiser is in charge}.

\subsubsection{Without optimisations (\texttt{-O0})}

Every local lives on the stack at a fixed offset from the frame pointer.
You can put a breakpoint and inspect any variable at any moment. Code is
slow and bloated but predictable. This is what you want when stepping
through code in a debugger.

\subsubsection{With optimisations (\texttt{-O2}, \texttt{-Os}, \texttt{-O3})}

\begin{itemize}[leftmargin=*]
    \item Locals that fit in registers and have short live ranges are kept
    \emph{entirely in registers} --- they never touch RAM.
    \item Locals whose value is never used after assignment are \emph{eliminated}.
    \item Constants computed at compile time are propagated into uses
    (constant folding).
    \item Dead code (\texttt{if (0) \{ ... \}}) is removed.
    \item Loops may be unrolled or rewritten.
\end{itemize}

\begin{warning}
A consequence that bites every beginner: in optimised code, a variable
\texttt{x} may not exist at the moment you place a breakpoint. The debugger
shows ``optimized out''. Use \texttt{-Og} (debug-friendly optimisation) or
\texttt{-O0} while learning. Also: an optimiser may remove a delay loop
unless variables are \texttt{volatile} or you use a real timer.
\end{warning}

\subsection{RISC vs.\ CISC variations}

\paragraph{RISC (ARM Cortex-M, RISC-V, AVR).}
Load--store architecture: ALU operates only on registers. To use a memory
value you must \texttt{LOAD} it, operate, then \texttt{STORE}.
RISCs have many general-purpose registers (ARM: 13 of them in user mode).
$\Rightarrow$ Compilers keep many locals in registers. Function prologues
push the registers they intend to clobber; epilogues pop them.

\paragraph{CISC (x86, x86-64).}
Memory operands allowed in most instructions: \texttt{ADD [ebp-4], eax}
adds \texttt{eax} directly into the local at \texttt{[ebp-4]}.
x86-32 has only \textasciitilde 8 GP registers, so locals more often
spill to the stack. x86-64 doubled the register count, so modern x86-64
code looks more RISC-like in its allocation patterns.

\paragraph{Practical impact.}
On a Cortex-M, a leaf function with five \texttt{int} locals usually
generates \emph{zero stack usage}: everything in R0--R12. On a 32-bit x86
the same function might push four of them to the stack. This is why
embedded stack budgets are often surprisingly small.

%======================================================================
\section{Debug Hardware Unit in Microcontrollers}
%======================================================================

On a PC, the debugger (gdb) talks to the OS, which uses \texttt{ptrace} to
control the target process. On an MCU there is no OS, so the debugger must
talk to \emph{hardware logic inside the chip} that can halt the core, peek
at registers, set breakpoints, and resume.

\subsection{The pieces (ARM CoreSight family)}

\begin{description}[leftmargin=*, style=nextline]
    \item[DAP (Debug Access Port).] The bridge between an external probe
    (ST-Link, J-Link) and the chip's internal debug bus.
    \item[JTAG / SWD.] The wire-level protocols the probe uses.
    JTAG = 4--5 wires (TCK, TMS, TDI, TDO, optional TRST), industry classic.
    SWD = 2 wires (SWCLK, SWDIO), ARM-specific, saves pins --- this is the
    common one on modern STM32.
    \item[FPB (Flash Patch \& Breakpoint).] Hardware breakpoint comparators.
    Limited (typically 4--8). The CPU halts on instruction fetch from a
    matching address.
    \item[DWT (Data Watchpoint \& Trace).] Watchpoints (halt when an
    address is read/written), plus cycle counters, exception trace, PC sampling.
    \item[ITM (Instrumentation Trace Macrocell).] Lets the firmware emit
    \texttt{printf}-style messages over a single SWO pin to the debugger.
    Insanely useful, costs almost nothing.
    \item[ETM (Embedded Trace Macrocell).] Streams the full instruction
    trace off-chip. Expensive in pin count, used for hard-to-reproduce
    timing bugs.
    \item[MPU (Memory Protection Unit).] Not strictly a debug feature, but
    pairs with debug: catches accesses to regions you've marked off-limits.
\end{description}

\subsection{Hardware vs.\ software breakpoints}

\begin{itemize}[leftmargin=*]
    \item \textbf{Hardware breakpoints} use FPB comparators. They work even
    on code in Flash (which you cannot modify at runtime). Limited number.
    \item \textbf{Software breakpoints} are implemented by overwriting an
    instruction with \texttt{BKPT} (Cortex-M) or \texttt{INT3} (x86), and
    restoring it on hit. Unlimited, but require writable code memory ---
    so on Flash they fall back to hardware breakpoints.
\end{itemize}

\subsection{What you can do thanks to all this}

\begin{itemize}[leftmargin=*]
    \item Halt the CPU and inspect every register and any memory location.
    \item Single-step instructions or source lines.
    \item Set a breakpoint that triggers only when \texttt{rx\_head == 64}.
    \item Watch a global and halt the moment something writes to it
    (great for ``who is corrupting my buffer?'').
    \item Reflash the chip directly from the IDE.
    \item Stream \texttt{printf} over SWO without using a UART.
\end{itemize}

%======================================================================
\section{Dynamic Memory Allocation in the C Standard Library}
%======================================================================

\texttt{malloc()} and \texttt{free()} are part of the C library, not the
language. Their implementation is up to whoever shipped your libc.

\subsection{What \texttt{malloc} actually does}

\begin{enumerate}[leftmargin=*]
    \item Maintains an internal data structure tracking free chunks in a
    pool of memory called the \textbf{heap}.
    \item On request, walks the free list, picks (or splits) a chunk big
    enough, marks it allocated, and returns a pointer just past the
    bookkeeping header.
    \item \texttt{free()} returns the chunk to the free list and (if
    possible) coalesces it with adjacent free chunks.
    \item When the heap is exhausted, \texttt{malloc} asks the OS (or, in
    embedded, an underlying \texttt{\_sbrk}) for more memory.
\end{enumerate}

\subsection{Embedded reality: newlib and \texttt{\_sbrk}}

Newlib's \texttt{malloc} doesn't know your chip's memory map. It calls
\texttt{\_sbrk(size)}, a function \emph{you} (or the framework) provide,
that hands out memory from a region typically defined in the linker script
as \texttt{\_end} up to \texttt{\_estack - stack\_reserve}.

\begin{lstlisting}[style=C, caption={Typical \_sbrk for STM32 with newlib}]
extern char _end;          /* end of .bss, start of heap (from linker) */
extern char _estack;       /* top of RAM (initial stack pointer)        */
static char *heap_ptr = &_end;

void *_sbrk(int incr)
{
    char *prev = heap_ptr;
    if (heap_ptr + incr > (&_estack - 0x400)) { /* 1 KiB stack reserve */
        errno = ENOMEM;
        return (void *) -1;
    }
    heap_ptr += incr;
    return prev;
}
\end{lstlisting}

\subsection{Why we don't trust \texttt{malloc} in embedded}

\begin{itemize}[leftmargin=*]
    \item \textbf{Non-determinism.} Allocation time depends on the current
    free-list state. A real-time loop can't tolerate this.
    \item \textbf{Fragmentation.} After many allocate/free cycles, the heap
    is full of small holes. A new allocation may fail even though the total
    free memory is enough. There is no compaction (you can't move allocated
    blocks --- pointers would dangle).
    \item \textbf{No graceful failure mode.} If \texttt{malloc} returns
    \texttt{NULL} on an MCU, what do you do? On a PC you log and bail out.
    On a motor controller, ``log and bail'' may mean ``the motor keeps
    spinning while the firmware is dead''.
    \item \textbf{Heap/stack collision.} Heap grows up, stack grows down,
    they share the same RAM. Worst case: heap silently overwrites stack
    or vice versa --- HardFault, far from the actual cause.
    \item \textbf{Code size.} Even if you ``never call malloc'', linking
    one function that calls \texttt{printf("\%f", ...)} can drag in malloc
    and float support and balloon your binary.
\end{itemize}

\begin{keyidea}
In safety-critical embedded code (automotive MISRA, aerospace DO-178C),
dynamic allocation is either entirely forbidden, or restricted to the
initialisation phase only.
\end{keyidea}

%======================================================================
\section{C Constructs to Avoid in Embedded}
%======================================================================

These rules look like nitpicks. They aren't. They are scar tissue from
real, expensive failures.

\subsection{Always use braces, even for single-statement blocks}

\begin{lstlisting}[style=C]
/* BAD */
if (flag)
    do_thing();
    log_event();           /* always runs! indentation lies */

/* GOOD */
if (flag) {
    do_thing();
}
log_event();
\end{lstlisting}

This is exactly the bug behind Apple's \emph{goto fail} TLS vulnerability
(2014): a duplicated \texttt{goto fail;} not guarded by braces skipped a
signature check unconditionally. Braces would have prevented it.
The MISRA-C 15.6 rule mandates braces.

\subsection{No unlimited recursion}

Recursion eats stack. The stack is small (often 1--8 KB) and shared with
ISRs. A recursion depth that depends on input data can silently overflow.
Either rewrite iteratively, or impose a hard cap and \emph{verify} that
the cap times the per-frame size fits in your stack budget.

\subsection{No \texttt{malloc}/\texttt{free}}

See previous section. Use one of:
\begin{itemize}[leftmargin=*]
    \item \textbf{Static allocation.} All buffers and objects declared at
    file scope.
    \item \textbf{Memory pools.} A fixed array of fixed-size blocks plus a
    bitmask of which are in use. \texttt{O(1)} alloc/free, no fragmentation.
    \item \textbf{Allocator-bounded design.} Allocate everything during
    \texttt{init()}, never afterwards. Then the system either boots or
    fails to boot --- never runs out of memory at runtime.
\end{itemize}

\subsection{Prefer indexes to pointers for array traversal}

\begin{lstlisting}[style=C]
/* Less safe */
for (int *p = arr; p < arr + N; ++p) *p = 0;

/* Safer, easier to bounds-check, easier to reason about */
for (size_t i = 0; i < N; ++i) arr[i] = 0;
\end{lstlisting}

Pointers do not carry array length: once you have a \texttt{int *p},
all the size information has evaporated. Indexes keep the array name
visible at every use, which makes bounds checks both possible and
obvious (\texttt{if (i < N)}). Tools like static analysers can also reason
much better about indices.

\subsection{Other classics worth knowing}

\begin{itemize}[leftmargin=*]
    \item No implicit \texttt{int} conversions / no shadowed variables.
    \item No \texttt{goto} except for ``goto cleanup'' patterns.
    \item No variable-length arrays (\texttt{int x[n]} where \texttt{n} is
    runtime --- they're stack-based and unbounded).
    \item No mixing of signed and unsigned in comparisons.
    \item Initialise every variable before use; don't trust ``zero by default''.
\end{itemize}

%======================================================================
\section{Dedicated Memory Areas for I/O Buffers and DMA Access}
\label{sec:dma}
%======================================================================

\textbf{DMA (Direct Memory Access)} is a hardware engine that copies data
between memory and a peripheral (or memory and memory) \emph{without
involving the CPU}. The CPU configures source, destination, size, and
goes back to doing other things; the DMA controller does the transfer
and raises an interrupt when done.

DMA is great for performance (megabytes/s) and power (CPU can sleep), but
it introduces constraints that ordinary code does not have.

\subsection{Why DMA needs special memory placement}

\begin{itemize}[leftmargin=*]
    \item \textbf{Bus accessibility.} The DMA controller has a specific
    matrix of which masters can reach which memory regions. On STM32F4,
    for instance, DMA2 can access SRAM1 but not CCM (Core-Coupled Memory).
    Putting a DMA buffer in CCM yields silent garbage.
    \item \textbf{Cache coherency.} On MCUs with a data cache (Cortex-M7),
    the CPU sees the cached copy and DMA sees the underlying RAM. If you
    write through the cache and the DMA reads RAM, it sees stale data; if
    the DMA writes RAM and the CPU reads the cache, it sees stale data.
    Solutions: place DMA buffers in a non-cacheable region, or
    invalidate/clean cache lines manually.
    \item \textbf{Alignment.} Many DMA engines require source/destination
    alignment to the transfer width (\texttt{\_\_attribute\_\_((aligned(4)))}).
    \item \textbf{Linker placement.} You usually carve out a specific
    region in the linker script (e.g.\ \texttt{SRAM2}) and force buffers
    into it via \texttt{\_\_attribute\_\_((section(".dma\_buffer")))}.
\end{itemize}

\subsection{An example layout}

\begin{lstlisting}[style=C]
/* aligned to 32 bytes (cache line), in non-cacheable SRAM2 */
__attribute__((aligned(32), section(".dma_buffer")))
static uint8_t uart_dma_tx[256];

__attribute__((aligned(32), section(".dma_buffer")))
static uint8_t uart_dma_rx[256];
\end{lstlisting}

%======================================================================
\section{The \texttt{const} and \texttt{volatile} Type Modifiers}
%======================================================================

These two modifiers terrify beginners and yet are absolutely fundamental
in embedded. They are independent: a variable can be \texttt{const},
\texttt{volatile}, both, or neither.

\subsection{\texttt{const}: ``the program promises not to write through this''}

\begin{itemize}[leftmargin=*]
    \item Tells the compiler the value will not change \emph{through this
    name}. Allows the compiler to put it in read-only memory (Flash) and to
    propagate the value into uses.
    \item For pointers, position matters:
        \begin{itemize}
            \item \texttt{const int *p}: pointer to const int. \texttt{*p}
            is read-only, \texttt{p} can be reassigned.
            \item \texttt{int * const p}: const pointer to int. \texttt{p}
            cannot be reassigned, \texttt{*p} is mutable.
            \item \texttt{const int * const p}: both.
        \end{itemize}
    \item Read it right-to-left: ``\texttt{p} is a const pointer to a const int''.
    \item \texttt{const} lookup tables in \texttt{.rodata} save precious RAM.
\end{itemize}

\subsection{\texttt{volatile}: ``the value may change without the compiler knowing''}

The compiler is allowed to assume that if you don't write to a variable,
its value doesn't change between reads. \texttt{volatile} explicitly
revokes that assumption. The compiler must emit an actual load from memory
every time you read it, and an actual store every time you write it. It
can't cache it in a register, can't elide ``redundant'' reads, can't reorder
across it freely.

When you need it:
\begin{enumerate}[leftmargin=*]
    \item \textbf{Memory-mapped hardware registers.} The peripheral changes
    them behind your back.
    \item \textbf{Variables shared with an ISR.} The ISR may write between
    two reads in main code.
    \item \textbf{Variables shared between cores / DMA / another bus master.}
    \item \textbf{Variables modified by \texttt{setjmp}/\texttt{longjmp}
    or signal handlers.}
\end{enumerate}

\subsection{\texttt{const volatile} --- yes, really}

A read-only status register: you can read it (compiler must do real load
every time), but you must not write it. Example: a hardware ``revision ID''
register.
\begin{lstlisting}[style=C]
#define CHIP_REV (*(const volatile uint32_t *)0x1FFF7A10)
\end{lstlisting}

\subsection{Famous \texttt{volatile} mistakes}

\begin{lstlisting}[style=C]
/* WRONG -- compiler may optimise this to an infinite check that never sees the flag change */
int flag = 0;
void ISR(void) { flag = 1; }
void wait(void) { while (!flag); }

/* RIGHT */
volatile int flag = 0;
\end{lstlisting}

\begin{warning}
\texttt{volatile} does \emph{not} make accesses atomic. Reading a 32-bit
\texttt{volatile} on a 32-bit MCU is one instruction (atomic), but a 64-bit
volatile is two stores. Sharing data between an ISR and main code wider
than the bus width still requires a critical section.
\end{warning}

%======================================================================
\section{Mandatory Content of \texttt{.c} and \texttt{.h} Files}
%======================================================================

C's compilation model is file-by-file. Each \texttt{.c} (\emph{translation
unit}) is compiled independently and must therefore see declarations of
everything it references from elsewhere. \texttt{.h} files are how you share
those declarations.

\subsection{Header (\texttt{.h}) file content}

\begin{itemize}[leftmargin=*]
    \item \textbf{Include guard} (or \texttt{\#pragma once}).
    \item \texttt{\#include}s for types this header uses (and \emph{only}
    those --- minimise transitive bloat).
    \item \texttt{typedef}s, \texttt{enum}s, and \texttt{struct} declarations
    that callers need to know about.
    \item \texttt{extern} declarations for globals exposed by the module.
    \item Function \textbf{prototypes} for the module's public API.
    \item Macros and inline functions that are part of the API.
    \item \emph{No definitions}: no function bodies (except \texttt{static
    inline}), no variable storage. A definition in a header pulled into
    multiple \texttt{.c} files yields multiple-definition link errors.
\end{itemize}

\begin{lstlisting}[style=C, caption={A clean header skeleton}]
#ifndef SENSOR_TEMP_H
#define SENSOR_TEMP_H

#include <stdint.h>
#include <stdbool.h>

typedef struct {
    float celsius;
    uint32_t timestamp_ms;
} temp_sample_t;

bool   temp_init(void);
bool   temp_read(temp_sample_t *out);
float  temp_to_kelvin(float c);

#endif /* SENSOR_TEMP_H */
\end{lstlisting}

\subsection{Source (\texttt{.c}) file content}

\begin{itemize}[leftmargin=*]
    \item \texttt{\#include "own\_header.h"} (first, to catch missing includes).
    \item \texttt{\#include} other headers as needed.
    \item \textbf{Static} (file-local) helper functions and variables.
    \item Definitions of the public functions declared in the header.
    \item No \texttt{extern} unless absolutely necessary --- if you need
    something external, declare it in a header.
\end{itemize}

\subsection{Function prototyping}

A \textbf{prototype} declares a function's name, return type, and parameter
types without a body:
\begin{lstlisting}[style=C]
int compute(int x, int y);   /* prototype */
\end{lstlisting}

Why it is mandatory in good embedded code:
\begin{itemize}[leftmargin=*]
    \item Without a prototype, the compiler assumes \texttt{int} return type
    and applies the old K\&R promotion rules to arguments. Pass a
    \texttt{float} to an unprototyped function and it's silently broken.
    \item It catches type mismatches at the call site.
    \item Modern C demands prototypes (\texttt{-Wmissing-prototypes},
    \texttt{-Wstrict-prototypes}) and treats \texttt{f()} differently from
    \texttt{f(void)} in C (the former still means ``unspecified args'').
\end{itemize}

\begin{note}
Rule of thumb: \emph{every} non-\texttt{static} function in a \texttt{.c}
file has a matching prototype in a \texttt{.h} file. Every \texttt{static}
function has a prototype near the top of its \texttt{.c} file so functions
can call each other freely without forward-declaration headaches.
\end{note}

%======================================================================
\section{Peripheral Representation in the Memory Map}
%======================================================================

In a microcontroller, peripherals (UART, SPI, GPIO, TIM, ADC, ...) are not
accessed via I/O instructions like on x86. They are \textbf{memory-mapped}:
each peripheral has a block of addresses that, when read or written by the
CPU, actually interact with the hardware.

\subsection{Three flavours of register}

Almost every peripheral exposes three categories:

\begin{description}[leftmargin=*, style=nextline]
    \item[Data register (DR).] Where the payload lives. Reading the UART
    DR yields the received byte. Writing it queues a byte to transmit.
    \item[Status register (SR).] Read-only flags reporting what the
    peripheral is doing: ``TX buffer empty'', ``RX byte available'',
    ``overrun error'', ``transmission complete''. You consult these to
    decide when to act.
    \item[Control register (CR).] Read/write bits that configure the
    peripheral: enable, baud rate, parity, interrupt enables, DMA enables,
    polarity, etc.
\end{description}

Larger peripherals add more (BRR for baud, GTPR for guard time, etc.) but
the DR/SR/CR pattern is the spine.

\subsection{Accessing registers from C: the typedef-struct trick}

Vendor headers (e.g.\ \texttt{stm32f4xx.h}) define each peripheral as a
\texttt{typedef}'d struct whose field offsets match the hardware layout.
A \texttt{\#define} maps a pointer at the peripheral's base address.

\begin{lstlisting}[style=C, caption={Simplified USART definition}]
typedef struct {
    __IO uint32_t SR;    /* 0x00: Status register      */
    __IO uint32_t DR;    /* 0x04: Data register        */
    __IO uint32_t BRR;   /* 0x08: Baud rate register   */
    __IO uint32_t CR1;   /* 0x0C: Control register 1   */
    __IO uint32_t CR2;   /* 0x10: Control register 2   */
    __IO uint32_t CR3;   /* 0x14: Control register 3   */
    __IO uint32_t GTPR;  /* 0x18: Guard time/prescaler */
} USART_TypeDef;

#define USART2 ((USART_TypeDef *) 0x40004400)

/* __IO expands to "volatile" -- absolutely critical here */
\end{lstlisting}

Using it is then natural C:
\begin{lstlisting}[style=C]
while (!(USART2->SR & USART_SR_TXE)) { /* wait */ }
USART2->DR = byte;
\end{lstlisting}

\begin{keyidea}
The \texttt{\_\_IO}/\texttt{volatile} qualifier is non-negotiable on
peripheral registers. Without it, the compiler will optimise away your
poll loop or fold multiple writes into one --- and the chip will behave
incomprehensibly.
\end{keyidea}

\subsection{Bit fields you'll meet}

\begin{itemize}[leftmargin=*]
    \item \textbf{r/w bits}: standard read/write.
    \item \textbf{r-only bits}: writes ignored. Status bits often.
    \item \textbf{w-only bits}: reads undefined. Control sequences sometimes.
    \item \textbf{rc\_w0 / rc\_w1 (read, clear by writing 0/1)}: a flag is set
    by hardware; you clear it by writing the specified value back. Very
    common for interrupt flags.
    \item \textbf{rs (set by reading)}: reading a register can change state
    (e.g.\ reading UART DR clears RXNE). This is why ``just looking at it''
    in the debugger can change behaviour.
\end{itemize}

%======================================================================
\section{STM32CubeIDE}
%======================================================================

STM32CubeIDE is ST Microelectronics' official Eclipse-based IDE for their
STM32 family. It is the most common environment in coursework because it
bundles everything you need and is free.

\subsection{What's inside}

\begin{itemize}[leftmargin=*]
    \item \textbf{Editor / build system.} Standard Eclipse CDT with managed
    Makefile generation.
    \item \textbf{Toolchain.} The \texttt{arm-none-eabi-} GCC family,
    pre-configured.
    \item \textbf{CubeMX (integrated).} A graphical configurator: you pick
    your chip, drag pins to peripherals, set clock trees, enable middleware
    (FreeRTOS, LwIP, FatFS, USB), and it \emph{generates} the initialisation
    C code, the linker script, and the startup file.
    \item \textbf{Debugger.} GDB front-end, talks to ST-Link / J-Link probes,
    supports breakpoints, watches, peripheral register view, RTOS-aware
    task list.
    \item \textbf{HAL/LL libraries.} Shipped as source so you can read and
    modify them.
\end{itemize}

\subsection{Typical workflow}

\begin{enumerate}[leftmargin=*]
    \item \textbf{New project} $\rightarrow$ pick board or MCU.
    \item In the \emph{Pinout \& Configuration} view, enable peripherals
    (e.g.\ ``USART2 in Asynchronous mode'', ``TIM3 in PWM Generation on CH1'').
    \item In \emph{Clock Configuration}, set your core, AHB, APB1, APB2
    clocks. Mismatches here are the source of many ``my UART speaks
    gibberish'' bugs.
    \item Click \emph{Generate Code}. The IDE produces \texttt{main.c}
    (with markers like \texttt{USER CODE BEGIN}/\texttt{END} you must
    write inside, otherwise re-generation overwrites your edits),
    \texttt{stm32xxxx\_hal\_msp.c}, the startup, and the linker script.
    \item Add your application code inside the user markers.
    \item Build. Connect probe. \emph{Debug} button flashes and starts
    a debug session.
\end{enumerate}

\subsection{FreeRTOS integration}

Inside CubeMX, enable \emph{Middleware $\rightarrow$ FREERTOS}. You get:
\begin{itemize}[leftmargin=*]
    \item Pre-configured \texttt{FreeRTOSConfig.h}.
    \item CMSIS-OS wrapper around FreeRTOS for portability.
    \item GUI to declare tasks, queues, mutexes, timers; CubeMX emits the
    \texttt{osThreadNew()} etc.\ calls.
    \item Automatic re-assignment of SysTick to FreeRTOS (and a different
    timer to HAL's tick source --- a classic pitfall).
\end{itemize}

%======================================================================
\section{HAL vs.\ LL Libraries; Callbacks; Weak Symbols}
%======================================================================

ST ships two parallel driver libraries for STM32 peripherals. Knowing
which is which and when to use which is half the battle.

\subsection{HAL (Hardware Abstraction Layer)}

\begin{itemize}[leftmargin=*]
    \item High-level, ``opinionated'' API: \texttt{HAL\_UART\_Transmit(\&huart2,
    buf, len, timeout)}.
    \item Hides peripheral state behind ``handles'' (\texttt{UART\_HandleTypeDef}).
    \item Same API across the whole STM32 family --- code portability between
    chips.
    \item Provides blocking, interrupt, and DMA variants of every operation.
    \item \textbf{Drawbacks:} larger code size, more cycles per call,
    sometimes opaque (you can't easily tell which registers it touches),
    the state machines occasionally have surprises (e.g.\ ``HAL\_BUSY''
    returned when you didn't expect).
\end{itemize}

\subsection{LL (Low-Level)}

\begin{itemize}[leftmargin=*]
    \item Thin, inline wrappers over register accesses:
    \texttt{LL\_USART\_TransmitData8(USART2, byte);}
    \item No state machines; almost no overhead --- often compiles to the
    same instructions as direct register code.
    \item Header-only, can be mixed with HAL on the same project.
    \item \textbf{Drawbacks:} you handle state yourself; less portability
    across families; less ``one-liner'' style operations.
\end{itemize}

\begin{note}
Common pragmatic choice: HAL for \emph{configuration} (cleaner, runs once
at boot), LL or raw registers for \emph{hot paths} (ISRs, tight loops).
\end{note}

\subsection{Interrupt handling through HAL/LL}

\textbf{With HAL.} You enable the IRQ in CubeMX. The IRQ handler in
\texttt{stm32xxxx\_it.c} calls e.g.\ \texttt{HAL\_UART\_IRQHandler(\&huart2)},
which detects which event fired, processes it, and then calls a
\emph{callback} that you override:
\begin{lstlisting}[style=C]
/* In your code: */
void HAL_UART_RxCpltCallback(UART_HandleTypeDef *h)
{
    if (h->Instance == USART2) {
        rx_byte_ready = 1;
    }
}
\end{lstlisting}

\textbf{With LL.} You write the IRQ handler directly:
\begin{lstlisting}[style=C]
void USART2_IRQHandler(void)
{
    if (LL_USART_IsActiveFlag_RXNE(USART2)) {
        rx_buf[rx_head++] = LL_USART_ReceiveData8(USART2);
    }
}
\end{lstlisting}

\subsection{Callbacks: the concept}

A \textbf{callback} is a function pointer that some library code will call
when an event occurs. The library defines \emph{when}, you define \emph{what}.
HAL uses callbacks pervasively for ISR events --- it provides the dispatch,
you write the handler.

\subsection{Weak symbols}

ST declares the default callback as ``weak'':
\begin{lstlisting}[style=C]
__attribute__((weak)) void HAL_UART_RxCpltCallback(UART_HandleTypeDef *h)
{
    UNUSED(h);
    /* default: do nothing */
}
\end{lstlisting}

A \textbf{weak symbol} is overridable: if any other translation unit defines
the same symbol \emph{without} \texttt{weak}, the linker uses the strong one
and silently ignores the weak one. If nobody defines it, the weak default
is used. This is how default ISR handlers and default callbacks work:
\begin{itemize}[leftmargin=*]
    \item Vendor provides a weak \texttt{Default\_Handler} as the address of
    every ISR vector slot. Spurious interrupts thus land in a sensible (if
    boring) infinite loop instead of executing garbage.
    \item You define a strong \texttt{TIM3\_IRQHandler} and the linker
    transparently replaces the weak default.
\end{itemize}

\begin{keyidea}
Weak symbols let a library ship sensible defaults while letting any user
project override them without modifying library files. They are the C
equivalent of ``virtual methods with a default implementation''.
\end{keyidea}

%======================================================================
\section{Sensor Interfacing: NTC Thermistor Example}
%======================================================================

A sensor produces a physical quantity (temperature, light, pressure) that
must be converted to a digital number. For a passive \textbf{analog sensor}
like an NTC thermistor, the path is:
\[
\text{Temperature} \longrightarrow \text{Resistance}
\longrightarrow \text{Voltage} \longrightarrow \text{ADC code}
\longrightarrow \text{Temperature (in software)}
\]

\subsection{NTC: what it is}

An NTC (\emph{Negative Temperature Coefficient}) thermistor is a resistor
whose resistance decreases as temperature increases. A 10 k$\Omega$ NTC has
$R = 10\,\mathrm{k}\Omega$ at 25\,$^\circ$C, perhaps 32\,k$\Omega$ at
0\,$^\circ$C and 4\,k$\Omega$ at 60\,$^\circ$C. The relationship is
strongly nonlinear, usually approximated by the \textbf{Beta equation}:
\[
\frac{1}{T} = \frac{1}{T_0} + \frac{1}{\beta}\ln\!\left(\frac{R}{R_0}\right)
\]
where $T$ is in Kelvin, $T_0 = 298.15$\,K, $R_0$ is the resistance at
$T_0$, and $\beta$ (Beta) is a material constant in the datasheet
(typically 3000--4500\,K).

A more accurate (and more compute-hungry) model is the \emph{Steinhart--Hart}
equation:
\[
\frac{1}{T} = A + B\ln(R) + C(\ln R)^3.
\]

\subsection{Resistance-to-voltage: the voltage divider}

The ADC reads voltage, not resistance. The simplest converter is a
\textbf{voltage divider}: NTC in series with a known reference resistor
$R_f$, the junction feeding the ADC.

\begin{center}
\begin{tabular}{c}
\texttt{Vcc ---[ R\_f ]---+---[ NTC ]--- GND} \\
\texttt{\hspace{36mm}|} \\
\texttt{\hspace{36mm}---> ADC input (V\_adc)}
\end{tabular}
\end{center}

By the divider rule:
\[
V_{\text{adc}} = V_{cc} \cdot \frac{R_{NTC}}{R_f + R_{NTC}}
\quad \Rightarrow \quad
R_{NTC} = R_f \cdot \frac{V_{\text{adc}}}{V_{cc} - V_{\text{adc}}}.
\]

The ADC then produces a code $D = \frac{V_{\text{adc}}}{V_{\text{ref}}}
\cdot (2^N - 1)$ where $N$ is the resolution. If $V_{\text{ref}} = V_{cc}$
they cancel out:
\[
R_{NTC} = R_f \cdot \frac{D}{(2^N - 1) - D}.
\]

\subsection{Design choices}

\begin{itemize}[leftmargin=*]
    \item Choose $R_f$ near $R_0$ to centre the divider's sensitive region
    around your temperature of interest.
    \item Watch self-heating: current through the NTC heats it. Keep the
    bias current small ($<0.1$ mA usually).
    \item ADC noise: take multiple samples and average; or use the MCU's
    oversampling feature.
    \item Decoupling: a small capacitor (e.g.\ 10\,nF) across the NTC
    smooths the input to the ADC sample-and-hold.
\end{itemize}

%======================================================================
\section{Fixed-Point vs.\ Floating-Point}
%======================================================================

How you represent fractional numbers has \emph{huge} performance and
accuracy implications on small MCUs.

\subsection{Floating-point: IEEE 754}

A float is sign $\cdot$ mantissa $\cdot 2^{\text{exponent}}$, encoded in
32 bits (\texttt{float}) or 64 bits (\texttt{double}).
\begin{itemize}[leftmargin=*]
    \item Huge dynamic range ($\approx 10^{-38}$ to $10^{+38}$ for float).
    \item Relative precision: ~7 decimal digits for float, ~15 for double.
    \item Arithmetic is non-associative (\texttt{(a+b)+c} $\neq$
    \texttt{a+(b+c)} in general).
\end{itemize}

\textbf{Hardware support} (FPU): Cortex-M4F/M7/M33 have a single-precision
FPU. \texttt{float} ops are one or a few cycles. Double precision is
software-emulated on M4F (slow), one-cycle on M7 with double-precision FPU.

\textbf{Without an FPU} (Cortex-M0/M0+/M3, AVR), all float arithmetic is
implemented in software --- a single multiplication costs hundreds of cycles
and \emph{kilobytes} of code (\texttt{libgcc} routines).

\subsection{Fixed-point: integers with an implied scale}

Pick a scale factor (power of two for efficiency) and interpret integers
as multiples of it. \textbf{Q-format} names them: \texttt{Q16.16} means
32 bits, 16 integer bits, 16 fractional bits, so the LSB represents
$2^{-16}$.

\begin{itemize}[leftmargin=*]
    \item Addition/subtraction: plain integer add/sub.
    \item Multiplication of \texttt{Qm.n} by \texttt{Qm.n} gives
    \texttt{Q(2m).(2n)} --- you must shift right by $n$ to renormalise
    (and use 64-bit intermediate to avoid overflow).
    \item Division: shift left by $n$ first, then integer divide.
\end{itemize}

\subsection{When to use which}

\begin{itemize}[leftmargin=*]
    \item \textbf{Cortex-M0/M3, AVR, no FPU}: fixed-point almost always.
    \item \textbf{Cortex-M4F+}: single-precision float is fine for most
    things and easier to reason about. Stick to \texttt{float}, not
    \texttt{double}, to keep using the FPU.
    \item \textbf{DSP / filtering at high rate}: fixed-point (often
    \texttt{Q15} or \texttt{Q31}) lets you use the SIMD/DSP instructions
    (Cortex-M4/M7 DSP extension, ARM CMSIS-DSP library).
    \item \textbf{Algorithms with wide dynamic range}: float, unless you
    really know your numeric range.
\end{itemize}

\begin{warning}
Beware of accidental promotions. \texttt{x / 2} with integer \texttt{x}
is integer division; \texttt{x / 2.0} promotes to double, drags in the
double-precision soft-float library, and bloats your binary. Write
\texttt{x / 2.0f} or just \texttt{x >> 1}.
\end{warning}

%======================================================================
\section{FreeRTOS Best Practices: Tasks, Queues, Events}
%======================================================================

\textbf{FreeRTOS} is a small preemptive RTOS. It gives you tasks (with
priorities), inter-task communication primitives, and synchronisation
primitives. The kernel is tiny (a few kB) and well understood.

\subsection{Tasks}

A \textbf{task} is a function with its own stack and priority that runs as
if it owns the CPU. The scheduler picks the highest-priority \emph{ready}
task and runs it; if a task blocks (waits on a queue, sleeps, waits on a
semaphore), the next-highest ready task runs.

Best practices:
\begin{itemize}[leftmargin=*]
    \item \textbf{Each task does one thing well.} ``UART RX task'',
    ``Motor control task'', ``Display update task''. Resist the temptation
    to make one giant task with a state machine.
    \item \textbf{Tasks block, they don't poll.} A loop with \texttt{vTaskDelay}
    is still polling. Block on the actual event (queue receive, semaphore
    take). The scheduler can then run other work or enter low-power.
    \item \textbf{Assign priorities by deadline criticality, not
    importance.} A 1-ms control loop is higher priority than the LCD
    redraw, even if the LCD ``feels'' more important.
    \item \textbf{Size stacks deliberately.} Use
    \texttt{uxTaskGetStackHighWaterMark} to measure actual usage.
    Allocate a bit of headroom but not megabytes.
    \item \textbf{Avoid \texttt{vTaskDelay} for timing of work.} Use
    \texttt{vTaskDelayUntil} for periodic tasks --- it compensates for the
    execution time of the loop body, giving you actual period accuracy.
\end{itemize}

\subsection{Queues}

A \textbf{queue} is a thread-safe FIFO. Producers \texttt{xQueueSend},
consumers \texttt{xQueueReceive}; both can block with timeouts.

\begin{itemize}[leftmargin=*]
    \item Pass copies of small structures through queues, not pointers to
    stack-local data.
    \item Use queues to convert ISR events into task work:
    \texttt{xQueueSendFromISR} in the ISR, \texttt{xQueueReceive} in the
    consumer task. The ISR stays short; the heavy work happens in task
    context.
    \item Use \texttt{portYIELD\_FROM\_ISR(xHigherPriorityTaskWoken)} at
    the end of the ISR so the consumer task can run immediately if it has
    higher priority than the interrupted task.
\end{itemize}

\subsection{Events and other primitives}

\begin{description}[leftmargin=*, style=nextline]
    \item[Binary semaphore.] One-bit signal. ``Something happened, go do it.''
    Released from ISR, taken by a task.
    \item[Counting semaphore.] Counts resources or events.
    \item[Mutex.] Mutual exclusion, with \emph{priority inheritance}
    (a low-prio task holding the mutex temporarily inherits the priority of
    a high-prio waiter, to avoid priority inversion).
    \item[Event group.] A 24-bit set of flags; a task can wait for ``any of
    these bits'' or ``all of these bits''. Good for ``wait for either
    serial-done or timeout''.
    \item[Direct task notifications.] The cheapest of all --- a per-task
    32-bit value with built-in semaphore semantics. Use them when only one
    task is waiting; they avoid allocating a queue/semaphore object.
\end{description}

\begin{warning}
Two notorious classes of RTOS bugs: \textbf{priority inversion} (use mutexes,
not binary semaphores, for resource protection) and \textbf{stack overflow}
(turn on \texttt{configCHECK\_FOR\_STACK\_OVERFLOW = 2} during development).
\end{warning}

%======================================================================
\section{Timer Peripheral: Prescaler and Period}
%======================================================================

A hardware timer is, at its core, a counter clocked from the system clock
(possibly divided). It counts from 0 up to a programmable maximum, then
either wraps or stops, and emits events (interrupt, DMA request, output
toggle) at programmable points.

\subsection{The two knobs}

\begin{description}[leftmargin=*, style=nextline]
    \item[Prescaler (PSC).] An integer divider applied to the timer's
    input clock before it reaches the counter. If the timer clock is
    $f_{\text{clk}}$ and the prescaler is $\text{PSC}$, the counter ticks
    at
    \[
    f_{\text{count}} = \frac{f_{\text{clk}}}{\text{PSC} + 1}.
    \]
    (The +1 is because the register value 0 means ``divide by 1''.)
    \item[Period / Auto-Reload Register (ARR).] The value at which the
    counter wraps (or the event fires). The update event happens every
    $\text{ARR}+1$ counter ticks, so the update frequency is
    \[
    f_{\text{update}} = \frac{f_{\text{clk}}}{(\text{PSC}+1)(\text{ARR}+1)}.
    \]
\end{description}

\subsection{Worked example}

Goal: produce a 1\,kHz tick from a 72\,MHz timer clock.
\[
\frac{72\,000\,000}{1\,000} = 72\,000 \text{ counter ticks per period}.
\]
A single 16-bit ARR maxes at 65535, so 72\,000 won't fit. Insert a small
prescaler:
\begin{itemize}[leftmargin=*]
    \item PSC = 71 $\Rightarrow$ counter ticks at $72\text{MHz}/72 = 1$ MHz.
    \item ARR = 999 $\Rightarrow$ update event every 1000 ticks
    $\Rightarrow$ 1\,kHz. Done.
\end{itemize}

\subsection{What you actually do with timers}

\begin{itemize}[leftmargin=*]
    \item \textbf{Periodic interrupt} (``tick''). Set PSC and ARR, enable
    update interrupt, run a handler at fixed cadence.
    \item \textbf{PWM output} (Output Compare mode). Counter compares
    against another register (CCR). Output flips at compare match,
    resets at ARR wrap $\Rightarrow$ duty cycle = CCR/(ARR+1).
    \item \textbf{Input capture}. Latch the counter value on an edge of
    an input pin --- measure pulse widths, frequencies, etc.
    \item \textbf{Encoder mode}. Two inputs decoded as quadrature: counter
    automatically tracks position of a rotary encoder.
    \item \textbf{Triggering ADC / DMA.} Timers can request a DMA transfer
    or start an ADC conversion at exact moments without CPU involvement.
    Foundation of high-rate, jitter-free signal acquisition.
\end{itemize}

%======================================================================
\section{Power Electronics from Scratch}
%======================================================================

Before the switching sections, the handful of device ideas they assume ---
built from the bottom up. Nothing new, just the floor.

\textbf{A transistor is a controlled valve.} A GPIO sources a few mA at
3.3\,V; motors/heaters/lamps want amps at 12--48\,V. You don't drive the
load from the pin --- you drive a transistor that drives the load. The
\emph{control terminal} (gate/base) is the handle the pin works for almost
nothing; the \emph{main path} (drain--source / collector--emitter) is the
pipe that carries the load current.

\textbf{Two kinds of valve.} A \emph{MOSFET} is \emph{voltage}-operated:
gate voltage (vs.\ its source) above $V_{th}$ opens the channel, and holding
it open costs $\approx$\,no current (a ``logic-level'' FET opens at 3.3\,V).
A \emph{BJT} is \emph{current}-operated: $I_C=\beta I_B$ in the active region
($\beta\approx$\,50--200), and it needs continuous base current to stay on.
That difference is why MOSFETs dominate switching and why a BJT's base drive
is a real design step (below).

\textbf{What turns a MOSFET on: gate-to-\emph{source}.} The channel opens
when $V_{GS}=V_G-V_S\ge V_{th}$ --- never gate-to-ground. So the whole
question is \emph{where the source sits}. Source at ground (low-side): a
3.3\,V pin is 3.3\,V of $V_{GS}$, trivial. Source up at Vcc (high-side): the
gate must go \emph{above} Vcc (a 12\,V rail wants $\sim$17\,V), which a logic
pin can't make alone.

\begin{note}
``It's gate-to-source, so where's the source?'' is the seed of the entire
low-/high-side story --- and the reason smart high-side drivers exist.
\end{note}

\textbf{BJT switch = overdriven into saturation.} An amplifier stays in the
active region; a \emph{switch} is driven hard until $V_{CE}$ bottoms out at
$V_{CE(\text{sat})}\approx0.1$--0.3\,V (saturation --- a closed switch with a
small drop). For margin you \emph{overdrive}: design to a low ``forced
$\beta$'' of 10--20 (feed $\sim I_C/15$) instead of the real
$\beta\approx100$. That ``forced $\beta$'' is what sizes the base resistor
(below).

\textbf{Making a voltage above the rail.} A high-side N-channel gate needs a
voltage above Vcc that isn't on the board --- so a \emph{charge pump} /
\emph{bootstrap} capacitor makes one. You rarely build it; it lives inside
integrated high-side switches.

\textbf{Why switching wastes energy.} Fully on, $V\approx0$; fully off,
$I\approx0$ --- so $P=VI\approx0$ in both steady states. Loss happens only
\emph{during the flip}, when $V$ and $I$ are briefly both present, giving the
$P_{\text{sw}}\approx\frac12 VI\,(t_r+t_f)\,f_{\text{sw}}$ you'll meet in the
next section. Faster edges $=$ less loss but more EMI.

\begin{note}
\textbf{Heat is Ohm's law renamed:} $\Delta T = P\cdot R_{th}$ (power
$\leftrightarrow$ current, $R_{th}\leftrightarrow$ resistance, temperature
rise $\leftrightarrow$ voltage), so $T_j=T_{\text{amb}}+P\,R_{th}$. A heatsink
is just a smaller $R_{th}$.
\end{note}

%======================================================================
\section{Load Switches: Low-Side and High-Side}
%======================================================================

A \textbf{load switch} is a transistor used to turn a load (motor, LED,
heater, relay coil) on and off. Where in the circuit you place the
switch fundamentally changes the design.

\subsection{Low-side switch}

The switch sits \emph{between the load and ground}. The load is
permanently connected to Vcc.

\begin{center}
\texttt{Vcc ---[ LOAD ]---+--- DRAIN/COLLECTOR \;\; SWITCH \;\; SOURCE/EMITTER ---GND}
\end{center}

To turn on:
\begin{itemize}[leftmargin=*]
    \item \textbf{NMOS as low-side}: drive the gate above the source (which is
    at ground), so $V_{GS} \approx V_{cc-logic}$. A 3.3\,V logic gate is
    plenty for a ``logic-level'' NMOS (e.g.\ AO3400, IRLML2502).
    \item \textbf{NPN BJT as low-side}: drive base current into the base
    resistor; $V_{BE} \approx 0.7$\,V; saturate the transistor.
\end{itemize}

\textbf{Advantages.} The control signal is referenced to ground, same as
the MCU. Easy to drive directly from a GPIO. Simple gate/base drive.

\textbf{Disadvantages.} The load is ``hot'' even when off (Vcc always
present at one terminal). If the load's downstream terminal accidentally
shorts to ground, the switch can't protect against it. Common-ground
EMI considerations.

\subsection{High-side switch}

The switch sits \emph{between Vcc and the load}. The load is permanently
connected to ground.

\begin{center}
\texttt{Vcc --- SOURCE/EMITTER \;\; SWITCH \;\; DRAIN/COLLECTOR ---+--- [ LOAD ] --- GND}
\end{center}

To turn on:
\begin{itemize}[leftmargin=*]
    \item \textbf{PMOS as high-side}: gate must be \emph{below} source by
    $|V_{GS(th)}|$. The source is at Vcc, so you must pull the gate down ---
    typically to GND. Driving a PMOS from a 3.3\,V MCU with Vcc=12\,V requires
    a level-shifter (an NPN/NMOS pulling the gate low).
    \item \textbf{NMOS as high-side}: source sits at the load, near Vcc when
    on. To keep it on you need $V_{GS} > V_{th}$, i.e.\ gate at $V_{cc} +
    V_{GS}$. That voltage doesn't exist on the board --- you must create it
    with a \emph{charge pump} or a \emph{bootstrap capacitor}.
\end{itemize}

\textbf{Advantages.} The load is fully isolated from Vcc when off ---
safer. A short-to-ground at the load side simply pulls more current
through the switch, which can be sensed and limited. Many industrial
applications mandate high-side switching for this reason.

\textbf{Disadvantages.} Drive circuitry is more complex. This is why
integrated high-side drivers exist (next section).

\subsection{NPN BJT as low-side switch: detailed walk-through}

An NPN turns on when $V_{BE} \gtrsim 0.6$\,V \emph{and} sufficient base
current flows. The relationship is:
\[
I_C = \beta \cdot I_B,
\]
but only in the active region. To use a BJT as a switch we want it
\emph{saturated} (very low $V_{CE}$, $\sim 0.2$\,V), which requires
\textbf{overdriving} the base, typically $I_B \approx I_C / 10$ (a
``forced beta'' of 10, well below the device's $\beta$ of perhaps 100).

%======================================================================
\section{Integrated High-Side Drivers}
%======================================================================

Because high-side NMOS driving is non-trivial, vendors sell single-chip
\textbf{high-side drivers}. Examples: BTS50xx (Infineon Profet), TPS272xx
(TI), VN5T family. Inside one tiny package you get:

\begin{itemize}[leftmargin=*]
    \item A power NMOS (or a clever PMOS stack) for the switch itself.
    \item A \textbf{charge pump} that generates a gate voltage above Vcc
    so the NMOS can be fully enhanced.
    \item \textbf{Current sense} (a small ``sense FET'' mirror or a low-side
    shunt) producing a proportional analog output (\texttt{IS} pin) for the
    MCU's ADC.
    \item \textbf{Current limit / over-current shutdown}, often configurable
    via an external resistor.
    \item \textbf{Thermal shutdown} when the die overheats.
    \item \textbf{Reverse-battery protection} in automotive variants.
    \item \textbf{Open-load detection} (in OFF state) for diagnostics.
    \item \textbf{Slew rate control}: the driver shapes the gate drive so
    that turn-on/turn-off is neither too slow (high $P = V \cdot I$ during
    transition $\Rightarrow$ heat) nor too fast (EMI, ringing on inductive
    loads).
\end{itemize}

\subsection{Why turn-on/turn-off speed matters}

\paragraph{Too slow.} The transistor spends time in its active region with
both high $V_{DS}$ and high $I_D$ --- power dissipation $P = V \cdot I$ peaks.
A 1\,A, 12\,V load with a 1\,$\mu$s transition wastes $\sim 6\,\mu$J per
edge. At 100\,kHz PWM that's 1.2\,W just from switching losses.

\paragraph{Too fast.} Sharp $dI/dt$ across parasitic inductance generates
voltage spikes that ring, radiate (EMI), and can damage the FET or trigger
spurious turn-ons of the opposite-side switch in a bridge.

\paragraph{Slew-rate-controlled drivers} compromise by ramping the gate
charging current, achieving a controlled (often programmable) $dV/dt$.

%======================================================================
\section{Switch Power Dissipation; NPN Base Resistor Sizing}
%======================================================================

\subsection{Where the power goes (MOSFET)}

\begin{itemize}[leftmargin=*]
    \item \textbf{Conduction loss}: $P_{\text{cond}} = I_D^2 \cdot R_{DS(on)}$.
    Buy a FET with low $R_{DS(on)}$ at the gate voltage you can supply.
    \item \textbf{Switching loss}: energy dissipated during transitions.
    Roughly $P_{\text{sw}} = \frac{1}{2} V_{DS} \cdot I_D \cdot (t_r + t_f)
    \cdot f_{\text{sw}}$.
    \item \textbf{Gate-drive loss}: $P_{\text{gate}} = Q_g \cdot V_{GS} \cdot
    f_{\text{sw}}$. Modest at low frequencies, can matter at MHz.
    \item \textbf{Body-diode loss}: in synchronous topologies, current may
    briefly flow through the body diode, which has a higher $V_F$ than the
    channel.
\end{itemize}

\subsection{Where the power goes (BJT)}

\begin{itemize}[leftmargin=*]
    \item \textbf{Conduction loss}: $P_{\text{cond}} = V_{CE(\text{sat})} \cdot I_C$.
    $V_{CE(\text{sat})}$ is typically 0.1--0.3\,V. Unlike a MOSFET this is
    a voltage drop, not an ohmic resistance: doubling the current does not
    quadruple the loss, only doubles it.
    \item \textbf{Base drive loss}: $P_{\text{base}} = V_{BE} \cdot I_B$
    (small but non-zero).
\end{itemize}

\subsection{Sizing an NPN base resistor}

Given:
\begin{itemize}[leftmargin=*]
    \item GPIO logic level: $V_{\text{IO}}$ (e.g.\ 3.3\,V).
    \item BJT $V_{BE}$ when saturated: $\approx 0.7$\,V.
    \item Desired collector current: $I_C$.
    \item Worst-case $\beta_{\min}$ from the datasheet (use this, not
    typical!).
    \item Overdrive factor (forced beta): pick $\beta_{\text{forced}} = 10$
    or so, to guarantee saturation across temperature/lot variation.
\end{itemize}

Required base current:
\[
I_B = \frac{I_C}{\beta_{\text{forced}}}.
\]
Voltage available across the resistor:
\[
V_R = V_{\text{IO}} - V_{BE}.
\]
Resistor value:
\[
R_B = \frac{V_R}{I_B} = \frac{V_{\text{IO}} - V_{BE}}{I_C / \beta_{\text{forced}}}.
\]

\textbf{Numerical example.} 3.3\,V GPIO, load 100\,mA, choose
$\beta_{\text{forced}} = 10$:
\[
I_B = 10\,\text{mA}, \quad V_R = 2.6\,\text{V}, \quad
R_B = 260\,\Omega.
\]
Use 220\,$\Omega$ (next standard E12 value down --- slight over-drive is
fine). Check the GPIO can source 10\,mA (most can, but at the edge of
spec).

\begin{warning}
\textbf{Never} drive a BJT base directly from a GPIO without a series
resistor --- you'll exceed the device's max base current and either kill
the BJT or kill the GPIO output stage.
\end{warning}

%======================================================================
\section{Full-Bridges and Shoot-Through}
%======================================================================

A \textbf{full-bridge (H-bridge)} is four switches arranged so a DC
voltage can be applied to a load in either polarity. The classical use:
driving a DC motor forwards and backwards.

\subsection{Topology}

\begin{center}
\begin{tabular}{cc}
\multicolumn{2}{c}{Vcc} \\
\multicolumn{2}{c}{|}\\
+---+---+ & +---+---+ \\
| Q1| & | Q3| \\
+---+ & +---+ \\
\multicolumn{2}{c}{ \quad MOTOR \quad } \\
+---+ & +---+ \\
| Q2| & | Q4| \\
+---+---+ & +---+---+ \\
\multicolumn{2}{c}{|} \\
\multicolumn{2}{c}{GND} \\
\end{tabular}
\end{center}

Q1 (high-side left) and Q4 (low-side right) on $\Rightarrow$ current flows
left-to-right through the load. Q3 + Q2 on $\Rightarrow$ current right-to-left.
PWM Q1 vs.\ Q2 (or Q3 vs.\ Q4) and you have direction + speed control.

\subsection{Shoot-through: the cardinal sin}

\textbf{Shoot-through} occurs when Q1 \emph{and} Q2 are both on
simultaneously (or Q3 and Q4), creating a direct Vcc-to-GND path through
two transistors. Currents of hundreds of amps can flow for microseconds.
The result: instant catastrophic failure --- molten silicon, flames,
all of the above.

This is rarely a steady-state problem (control logic won't intentionally
command it). It happens during \emph{transitions}: when you ask Q1 to
turn off and Q2 to turn on, the real-world devices don't switch
instantly. If Q1's turn-off is slower than Q2's turn-on, there's an
overlap window where both are partially on.

\subsection{Protection: dead-time}

The solution is \textbf{dead-time}: insert a short period where
\emph{both} switches in a leg are off, between turning one off and
turning the other on. During dead-time, current continues to flow through
the body diodes of the off switches (or through external freewheel
diodes). Typical dead-times: 100\,ns to a few \textmu s, depending on
device speed.

\subsubsection{Timer-based dead-time}

Advanced timers in modern MCUs (STM32 TIM1, TIM8) support
\textbf{complementary outputs with programmable dead-time} natively. You
configure one PWM channel and the timer automatically generates a
complement on a second pin, separated by your specified dead-time
(register DTG --- ``Dead-Time Generator''). No software fiddling, hardware
guarantees the gap.

\subsubsection{Driver-based dead-time}

Half-bridge gate driver chips (e.g.\ IR2104, DRV8301, UCC27210) include
built-in dead-time insertion. You provide a single PWM signal; the
driver generates the high-side and low-side gate signals with an
internal interlock that physically prevents both from being asserted
simultaneously.

\begin{keyidea}
Belt-and-braces is normal: \emph{both} the timer and the driver enforce
dead-time. A software bug that disables one is then still caught by the
other.
\end{keyidea}

%======================================================================
\section{Digital-Interface Sensors: The Accelerometer Example}
%======================================================================

A \textbf{digital-interface sensor} performs the analog-to-digital
conversion internally and exposes its measurements via a digital bus
(I\textsuperscript{2}C, SPI, sometimes UART or a 1-wire variant). The MCU
just reads registers.

\subsection{Why digital is nicer}

\begin{itemize}[leftmargin=*]
    \item No noisy analog wires snaking across the board.
    \item Calibration done once at the factory and baked into the sensor.
    \item Configurable filters, ranges, sample rates accessible via
    registers.
    \item Built-in features: temperature compensation, motion-detection
    interrupts, FIFOs.
\end{itemize}

\subsection{MEMS accelerometer, briefly}

An accelerometer measures linear acceleration on 1, 2, or 3 axes. A MEMS
device contains a tiny silicon proof mass on springs (etched into the
chip); acceleration deflects the mass, changing capacitances measured by
internal circuitry. Output: $g$-force values in the configured range
(e.g.\ $\pm 2g, \pm 4g, \pm 8g, \pm 16g$).

\subsection{Example: hard disk drives and free-fall detection}

Many laptop hard drives ship with a 3-axis accelerometer. Use case: if
the drive senses \emph{free fall} (all three axes near 0\,$g$
simultaneously, i.e.\ no gravity reaction --- only happens in free fall),
firmware concludes the drive is falling. Within milliseconds it parks the
head off the platter, so when the laptop hits the floor the platter and
head do not collide.

The detection is built into the sensor: registers configure a free-fall
threshold and duration; an interrupt pin asserts when the condition
holds. The MCU's ISR triggers the head-park sequence. The whole loop is
fast precisely because the heavy lifting is in the sensor.

Other classic uses: smartphones (rotate the display), pedometers,
vibration monitoring of industrial machinery, automotive airbag
triggering.

%======================================================================
\section{I\textsuperscript{2}C and SPI Protocols}
%======================================================================

These are the two ubiquitous embedded buses for chip-to-chip
communication on a board. Memorise them.

\subsection{I\textsuperscript{2}C (Inter-Integrated Circuit)}

A 2-wire, multi-master, multi-slave bus.

\begin{description}[leftmargin=*, style=nextline]
    \item[Wires] SDA (Serial Data, bidirectional, open-drain) and SCL
    (Serial Clock, open-drain). Pulled up to Vcc by external resistors
    (typically 4.7--10\,k$\Omega$).
    \item[Speeds] Standard 100\,kHz, Fast 400\,kHz, Fast+ 1\,MHz,
    High-Speed 3.4\,MHz.
    \item[Addressing] Each slave has a 7-bit address (some support 10-bit).
    A transaction starts with the master pulling SDA low while SCL is high
    (\emph{START condition}), then clocking out the 7-bit address + 1 R/W bit.
    The matching slave \texttt{ACK}s by pulling SDA low for one clock.
    \item[Direction] Set by the R/W bit. Reads typically require: send
    register address (write phase), \emph{repeated START}, then read.
    \item[Stop] Master releases SDA from low to high while SCL is high
    (\emph{STOP condition}).
\end{description}

\textbf{Advantages.} Only 2 wires for an entire bus of many devices.
Standard, supported everywhere. Software-emulatable on plain GPIO.

\textbf{Drawbacks.} Slower than SPI. Open-drain means rise time is limited
by the pull-up resistor and the bus capacitance. Bus can hang if a slave
holds SDA low (recovery: toggle SCL nine times). Addressing collisions
between vendors --- watch the address map of every chip you put on the bus.

\subsection{SPI (Serial Peripheral Interface)}

A 4-wire (sometimes 3-wire), full-duplex, master-slave bus.

\begin{description}[leftmargin=*, style=nextline]
    \item[Wires] SCK (clock, master $\rightarrow$ slave), MOSI (Master Out
    Slave In), MISO (Master In Slave Out), and one CS (Chip Select) per
    slave (active low). Push-pull, no pull-ups needed.
    \item[Speeds] Tens of MHz easily, up to 100+ MHz on capable devices.
    \item[Operation] Master asserts CS, then shifts bits out on MOSI on
    one clock edge while the slave shifts back on MISO on the other edge.
    Every clock = one bit each direction; full-duplex.
    \item[Modes] Four combinations of clock polarity (CPOL) and clock phase
    (CPHA): Mode 0 to Mode 3. The chosen mode must match between master and
    slave. Wrong mode = nothing works, often silently.
    \item[Daisy-chaining] Some devices can be chained: MISO of one goes to
    MOSI of the next. Useful for shift registers.
\end{description}

\textbf{Advantages.} Very fast, simple state machine, no addressing
overhead, push-pull drives don't care about pull-ups.

\textbf{Drawbacks.} One CS line per slave (more pins). No standard
framing or error detection --- it's just shift registers, you pay
attention to bytes. No flow control --- slave must be ready every clock.
No defined slave-to-master interrupt mechanism (often a separate GPIO
``INT'' pin handles this).

\subsection{Comparison cheat sheet}

\begin{tabular}{p{0.18\linewidth} p{0.39\linewidth} p{0.39\linewidth}}
\toprule
& \textbf{I\textsuperscript{2}C} & \textbf{SPI} \\
\midrule
Wires (n slaves) & 2 total & 3 + n (one CS each) \\
Speed & up to a few MHz & tens to 100+ MHz \\
Multi-master & Yes (arbitration) & No (one master typically) \\
Addressing & 7/10-bit on bus & Per-slave CS line \\
Duplex & Half (one direction at a time) & Full \\
Framing/ACK & Built-in & None (raw bits) \\
Drive type & Open-drain + pull-up & Push-pull \\
Typical use & Sensors, EEPROMs, small ICs & SD card, displays, ADCs, fast sensors \\
\bottomrule
\end{tabular}

\subsection{In an STM32 (peripherals view)}

Each MCU has hardware peripherals (e.g.\ \texttt{I2C1, I2C2, I2C3,
SPI1, SPI2, SPI3}) with their own DMA channels and interrupts.
Workflow:
\begin{enumerate}[leftmargin=*]
    \item In CubeMX, enable the peripheral, set pins, baud/speed, SPI
    mode, optional DMA.
    \item In code, use HAL or LL: \texttt{HAL\_I2C\_Mem\_Read(\&hi2c1,
    addr<<1, reg, 1, buf, n, timeout)} for I\textsuperscript{2}C reads.
    \item For high-rate transfers, use the \texttt{\_DMA} variants and
    handle the completion callback to chain the next transaction.
\end{enumerate}

\begin{note}
A subtle gotcha: HAL's \texttt{HAL\_I2C\_*} functions expect the slave
address \emph{shifted left by 1}, because the LSB is used internally for
the R/W bit. Forgetting the shift gives ``no ACK'' and confused debugging.
\end{note}

%======================================================================
\section*{Closing Thoughts}
%======================================================================
\addcontentsline{toc}{section}{Closing Thoughts}

This syllabus is broad on purpose: firmware design lives at the
intersection of \emph{C language}, \emph{computer architecture},
\emph{electronics}, and \emph{real-time systems}. The recurring themes:

\begin{itemize}[leftmargin=*]
    \item Embedded code runs against \emph{physics}, not a forgiving OS.
    Determinism, resource budgets, and correct hardware interaction
    matter more than elegance.
    \item Understanding what the compiler and linker actually do
    (sections, optimisation, calling conventions, weak symbols) is not
    optional --- it's the only way to debug ``the C is fine but the chip
    is wrong''.
    \item Hardware peripherals are just registers. Once you read the
    reference manual you can drive anything; libraries are conveniences,
    not magic.
    \item Patterns repeat: ISR pushes data to a queue, task processes it;
    timer triggers ADC via DMA; weak default ISR overridden by your
    strong one; mutex protects shared state with priority inheritance.
    Recognising these patterns lets you read any embedded code base.
\end{itemize}

Now go flash a board and break something. That's how this stuff actually
sticks.

\end{document}
